\begin{center}
\LARGE \textbf{Обходим Великий Российский Файрвол}\\
\large Часть 2 из 5\\
\normalsize\textbf{В России официально установлена цензура Интернета. Но контролировать Интернет не так просто, как это может показаться с первого взгляда.}
\end{center}
\begin{figure}[h]
\center{\includegraphics[width=0.25\linewidth]{Tor}}
\end{figure}
\textbf{Tor} --- анонимная оверлейная сеть, использующая принцип луковой маршрутизации, исходные коды которой распространяются на условиях лицензии BSD.\\
\textbf{Луковая маршрутизация} --- технология анонимного обмена информацией, использующая многократное шифрование и пересылку через цепочки узлов. Каждый луковый маршрутизатор в цепочке удаляет слой шифрования и пересылает сообщение дальше, согласно полученным инструкциям, где все повторится. И так до тех пор, пока сообщение не достигнет адресата. Такое название технология получила из-за сходства данного процесса с очисткой луковицы.
\begin{important}
Никогда не используйте анонимные сети с настройками, которые позволяют использовать Интернет в одном браузере в обход самой сети. В этом случае будет достаточно вставить любое изображение из внешнего Интернета, чтобы получить ваш реальный IP-адрес.
\end{important}
\subsection*{Установка}
Для установки посетите \url{https://torproject.org} или установите пакет с помощью пакетного менеджера вашего дистрибутива.
\subsection*{Использование}
\begin{important}
Не забывайте, что при использовании обычного Интернета через Tor, последняя нода в цепочке (exit-нода) видит трафик нешифрованным.
\end{important}
Самым простым способом использования Tor является установка Tor Browser Bundle. Просто скачайте его с \url{https://torproject.org/projects/torbrowser.html}, распакуйте и запустите.\\
Вы также можете установить Tor и задать в настройках приложений, которые вы хотите через него использовать, адрес socks5-прокси 127.0.0.1:9050.\\
Для использования приложений через Tor, не имеющих настроек прокси, скачайте torsocks --- \url{https://code.google.com/p/torsocks}.
\subsection*{Некоторые скрытые сервисы}
\begin{enumerate}
\item \url{http://dppmfxaacucguzpc.onion} --- TorDir, каталог скрытых сервисов.
\item \url{http://jhiwjjlqpyawmpjx.onion} --- TorMail, электронная почта в Tor.
\item \url{http://silkroadvb5piz3r.onion} --- Silk Road, анонимная торговая площадка, принимающая оплату в Bitcoin.
\item \url{http://4eiruntyxxbgfv7o.onion/pm} --- TorPM, сервис обмена сообщениями.
\item \url{http://c4wcxidkfhvmzhw6.onion} --- PrivacyBox, сервис анонимных контактных форм.
\end{enumerate}
\vfill
\scriptsize Создано на основе материалов из <<Настольной книги анонима>> --- \url{anonhandbook.org}, \url{anonhandbook.i2p}, \url{oxgzwiiypou6udlp.onion}
\normalsize
